% ==========================================================
%  OLD EXAM QUESTIONS — DL Architectures & Generative Techniques
%  Built from the 28.06.2024 main exam and the 2022 / 2024 retries.
%  Shares the visual language of summary.tex (article class).
% ==========================================================

\documentclass[11pt,oneside]{article}

\usepackage[
  a4paper,
  top=2.4cm,
  bottom=2.1cm,
  left=2.3cm,
  right=2.3cm,
  headheight=15pt
]{geometry}

\usepackage[T1]{fontenc}
\usepackage[utf8]{inputenc}
\usepackage{lmodern}
\usepackage{microtype}
\usepackage{inconsolata}
\usepackage{amsmath,amssymb,mathtools,bm}
\usepackage{enumitem}
\usepackage{booktabs,array,multicol,longtable}
\usepackage{xcolor}
\usepackage{graphicx}
\usepackage[most]{tcolorbox}
\usepackage{titlesec}
\usepackage{fancyhdr}
\usepackage[hidelinks]{hyperref}

\setlength{\parindent}{0pt}
\setlength{\parskip}{0.55em}
\setlength{\arraycolsep}{2.5pt}   % tighten matrices in answer options
\setlist[itemize]{topsep=3pt,itemsep=2pt,leftmargin=2.2em}
\setlist[enumerate]{topsep=3pt,itemsep=2pt,leftmargin=2.2em}

% ── Colors ───────────────────────────────────────────────
\definecolor{accent}{HTML}{1B3A5C}
\definecolor{q@frame}{HTML}{1A5276}
\definecolor{q@bg}{HTML}{F6FAFD}
\definecolor{key@frame}{HTML}{1E8449}
\definecolor{key@bg}{HTML}{F4FBF6}
\definecolor{note@frame}{HTML}{9A6A00}
\definecolor{note@bg}{HTML}{FEF9E7}

% ── Box global style ─────────────────────────────────────
\tcbset{
  enhanced,
  breakable,
  arc=2.5mm,
  boxrule=0.5pt,
  left=4.5mm,
  right=4mm,
  top=2.5mm,
  bottom=2.5mm,
  fonttitle=\bfseries\small,
  sharp corners=west,
  parbox=false
}

\newtcolorbox{qbox}[1][]{
  colback=q@bg,
  colframe=q@frame,
  borderline west={3.5pt}{0pt}{q@frame},
  title={#1}
}

\newtcolorbox{keybox}[1][]{
  colback=key@bg,
  colframe=key@frame,
  borderline west={3.5pt}{0pt}{key@frame},
  title={#1}
}

\newtcolorbox{notebox}[1][]{
  colback=note@bg,
  colframe=note@frame,
  borderline west={3.5pt}{0pt}{note@frame},
  title={#1}
}

% ── Section headings ─────────────────────────────────────
\titleformat{\section}
  {\normalfont\Large\bfseries\color{accent}}
  {}{0em}{}
  [{\vspace{1pt}\color{accent!45}\titlerule[0.45pt]}]
\titlespacing*{\section}{0pt}{3.2ex plus 1ex minus .2ex}{1.8ex plus .2ex}

\titleformat{\subsection}
  {\normalfont\normalsize\bfseries\color{accent!80!black}}
  {}{0em}{}
\titlespacing*{\subsection}{0pt}{2.2ex plus 1ex minus .2ex}{1.0ex plus .2ex}

% ── Header & footer ──────────────────────────────────────
\pagestyle{fancy}
\fancyhf{}
\fancyhead[L]{\small\color{gray!70}Old Exam Questions — \coursetitle}
\fancyhead[R]{\small\color{gray!70}\thepage}
\renewcommand{\headrulewidth}{0.4pt}

% ── Document metadata ─────────────────────────────────────
\newcommand{\coursetitle}{DL Architectures \& Generative Techniques}

% ── Question counter & macros ────────────────────────────
\newcounter{question}
\newcommand{\questiontitle}[1]{%
  \stepcounter{question}\begin{qbox}[Question \thequestion: #1]}
\newcommand{\closequestion}{\end{qbox}}

\newcommand{\repeatnote}{\textit{\color{gray!70}Asked multiple times.}}
\newcommand{\choice}{\item[$\square$] }
\newcommand{\answerline}[1]{\makebox[#1]{\hrulefill}}

\newcommand{\workingspace}[1]{%
  \par\vspace{2pt}%
  \begin{tcolorbox}[
    enhanced,
    breakable=false,
    colback=white,
    colframe=gray!40,
    boxrule=0.45pt,
    arc=0pt,
    sharp corners,
    height=#1,
    left=3mm, right=3mm, top=3mm, bottom=3mm,
    frame style={dashed}
  ]%
  \end{tcolorbox}%
  \vspace{2pt}%
}

\newcommand{\examfigure}[2][width=0.82\linewidth]{%
  \begin{center}
    \includegraphics[#1]{#2}
  \end{center}
}

% ── Math macros ──────────────────────────────────────────
\newcommand{\R}{\mathbb{R}}
\newcommand{\E}{\mathbb{E}}
\newcommand{\Prob}{\mathbb{P}}
\newcommand{\KL}[2]{D_{\mathrm{KL}}\!\left(#1\,\|\,#2\right)}
\newcommand{\dd}{\,\mathrm{d}}
\newcommand{\Cov}{\mathrm{Cov}}
\newcommand{\loss}{\mathcal{L}}
\newcommand{\Normal}{\mathcal{N}}
\newcommand{\given}{\,\vert\,}
\newcommand{\T}{^{\mathsf{T}}}
\DeclareMathOperator{\softmax}{softmax}


% ==========================================================
%  BEGIN DOCUMENT
% ==========================================================
\begin{document}

% ── Title page ───────────────────────────────────────────
\thispagestyle{empty}
\vspace*{1.4cm}
\begin{center}
  {\color{accent}\rule{\linewidth}{1.8pt}}\par
  \vspace{0.65cm}
  {\fontsize{26}{32}\selectfont\bfseries\color{accent}Old Exam Questions\\[2pt] \coursetitle\par}
  \vspace{0.55cm}
  {\color{accent!35}\rule{\linewidth}{0.5pt}}\par
  \vspace{0.55cm}
  {\large Stefan Eder\par}
  \vspace{0.15cm}
  {\small\color{gray!70}\today\par}
\end{center}

\vspace{0.9cm}
\begin{notebox}[How this file is organized]
Collected from the main exam (28.06.2024) and the 2022 / 2024 retry exams. Questions are
grouped by topic so that the sections line up with the chapters of the course summary.
Exact duplicates were removed; when essentially the same question recurred across exams it is
marked with \repeatnote\ inside the box. Several originals were figure-based (drag-and-drop
annotations, hand-drawn graphs, plots); where screenshots are available, they are included
directly. The mathematical content and answer choices are kept intact; the answer key is at the end.
\end{notebox}

\tableofcontents
\newpage


% ==========================================================
\section{Image Classification \& CNN Architectures}
% ==========================================================

\questiontitle{AlexNet}
Which of the following statements about AlexNet (Krizhevsky et al., 2012) are correct?
(Select one or more.)
\begin{itemize}[leftmargin=1.8em]
  \choice AlexNet does not have fully-connected layers.
  \choice AlexNet uses strided convolutions instead of max-pooling for down-sampling.
  \choice AlexNet is a deep convolutional network with more than 1 hidden layer.
  \choice One of the most important components of AlexNet is BatchNorm.
  \choice AlexNet uses input images with a resolution of $224\times224$ pixels and three channels.
\end{itemize}
\closequestion

\questiontitle{VGGNet}
Which of the following statements about VGGNet are correct? (Select one or more.)
\begin{itemize}[leftmargin=1.8em]
  \choice It predominantly uses $3\times3$ convolutions.
  \choice It is a single-stage approach (a single $5\times5$ layer rather than stacked $3\times3$ layers).
  \choice It has a symmetric encoder--decoder structure.
  \choice The majority of the network's weights are located in the fully-connected layers.
  \choice It cannot be used for image classification.
\end{itemize}
\closequestion

\questiontitle{Normalization in Image Classifiers}
Which of the following statements about normalization in image-classification networks are correct?
(Select one or more.)
\begin{itemize}[leftmargin=1.8em]
  \choice Batch Normalization has been used since AlexNet.
  \choice Since Xception, most image classifiers started to use Layer Normalization.
  \choice Batch Normalization was used in InceptionNet to speed up training.
  \choice Batch Normalization is part of most deep networks for image classification (e.g.\ ResNet).
\end{itemize}
\closequestion

\questiontitle{Residual Networks}
Which of the following statements about Residual Networks (ResNets) are correct? (Select one or more.)
\begin{itemize}[leftmargin=1.8em]
  \choice Residual networks use skip-connections that add the representation $\mathbf{h}^{l}$ of a
        former layer $l$ to a later layer $l+1$, i.e.\ $\mathbf{h}^{l+1}=\mathbf{h}^{l}+F(\mathbf{h}^{l})$,
        where $F$ can be any transformation with appropriate dimensions (e.g.\ a fully-connected layer).
  \choice Because of skip-connections, very deep ResNets can be trained even without normalization techniques.
  \choice Pre-trained ResNets cannot be used for transfer learning because the learned features are usually
        overly specified to the task on which they were trained.
  \choice Residual networks are always convolutional and cannot be used in purely fully-connected networks.
  \choice Leaky ReLUs cannot be used in Residual Networks.
\end{itemize}
\closequestion

\questiontitle{$1\times1$ Convolution Output Size}
You add a $1\times1$ convolutional layer with 16 kernels, padding 0, stride 1, to an input of
shape $[128,128,64]$. Which statement is correct?
\begin{itemize}[leftmargin=1.8em]
  \choice It is not possible to apply this convolution layer.
  \choice The resulting dimension is $[128,128,16]$.
  \choice A $1\times1$ convolution does not exist.
  \choice Better: replace the $1\times1$ layer with $2\times2$ max-pooling --- an equivalent but more efficient operation.
  \choice Better: use $1\times1$ with stride 2 --- equivalent but more efficient.
\end{itemize}
\closequestion

\questiontitle{Receptive Field of the Last Layer}
You are given a CNN with a fixed architecture. You may only change the parameters of the
\emph{last} convolutional layer, and you want to strongly increase the receptive field (RF)
of each neuron in that last layer. For each change, decide whether it produces
\textbf{no increase}, a \textbf{small increase}, or a \textbf{strong increase} of the RF.
\begin{itemize}[leftmargin=1.8em]
  \item Increase the stride of the convolution \answerline{4cm}
  \item Switch to dilated convolutions \answerline{4cm}
  \item Increase the number of kernels \answerline{4cm}
  \item Increase the kernel size \answerline{4cm}
\end{itemize}
\closequestion

\questiontitle{Zero Initialization --- Forward Pass}
You design a standard 2-layer fully-connected network with sigmoid activations and SGD.
You initialize all weights and biases to zero and forward-propagate an input $\mathbf{x}$.
What is the output $\hat y$?
\begin{itemize}[leftmargin=1.8em]
  \choice $0$
  \choice $e^{x}$
  \choice $0.5$
  \choice $-1$
  \choice $1$
\end{itemize}
\closequestion

\questiontitle{Zero Initialization --- Weight Update}
Consider the same all-zero-initialized network. You forward-propagate a batch, backpropagate,
and update the parameters once. $\mathbf{W}^{(2)}$ denotes the weight matrix of the second layer.
Which statement is then true?
\begin{itemize}[leftmargin=1.8em]
  \choice The entries of $\mathbf{W}^{(2)}$ are all positive.
  \choice The entries of $\mathbf{W}^{(2)}$ may be positive or negative.
  \choice The entries of $\mathbf{W}^{(2)}$ are all negative.
  \choice The entries of $\mathbf{W}^{(2)}$ are all zero.
\end{itemize}
\closequestion


% ==========================================================
\section{Semantic Segmentation}
% ==========================================================

\questiontitle{Self-Driving Scene Labelling}
A camera in a self-driving car provides road-scene images, and you train a network that
outputs an image in which every pixel is annotated as car, road, pedestrian, etc.
The figure shows examples of the input images and the corresponding pixel-wise labels.
\examfigure[width=0.88\linewidth]{figures/exam-semantic-segmentation-road-scene.png}
Which statements are true in this context? (Select one or more.)
\begin{itemize}[leftmargin=1.8em]
  \choice To train such a network, labels for each pixel of the training images are required.
  \choice This represents an object-detection task in computer vision.
  \choice Generative adversarial networks are the most appropriate type of network for this task.
  \choice Deep networks cannot be used because there is no appropriate loss function.
  \choice The intersection over union (IoU) would not be an appropriate metric to assess such a model.
\end{itemize}
\closequestion

\questiontitle{IoU Computation}
The figure shows a ground-truth object region (solid brown outline) and a predicted object region
(dashed green outline) on a grid of pixels. The ground-truth region covers $10$ pixels, the
predicted region covers $8$ pixels, and they overlap in $5$ pixels. Compute the IoU metric,
rounded to three decimals.
\examfigure[width=0.58\linewidth]{figures/exam-iou-pixel-regions.png}

Answer: \answerline{3cm}
\closequestion


% ==========================================================
\section{Object Detection \& Localization}
% ==========================================================

\questiontitle{Object-Detection Methods}
Which of the following statements about deep-learning-based object-detection methods are correct?
(Select one or more.)
\begin{itemize}[leftmargin=1.8em]
  \choice Faster R-CNN uses a region proposal network rather than the selective-search algorithm used in R-CNN.
  \choice The operations RoI-pooling and RoI-align operate on feature maps and not on the input image.
  \choice R-CNN uses AlexNet to extract features for Regions of Interest.
  \choice Object detection can be formulated as a regression task in a multi-task setting.
  \choice Fast R-CNN has two output layers (heads): one for regression and one for classification.
\end{itemize}
\closequestion

\questiontitle{One-Stage vs.\ Two-Stage Detectors}
For each architecture, state whether it is a \textbf{one-stage} or \textbf{two-stage} detector.
\begin{itemize}[leftmargin=1.8em]
  \item Faster R-CNN \answerline{3.5cm}
  \item YOLO \answerline{3.5cm}
  \item Mask R-CNN \answerline{3.5cm}
  \item Fast R-CNN \answerline{3.5cm}
  \item R-CNN \answerline{3.5cm}
  \item SSD \answerline{3.5cm}
\end{itemize}
\closequestion

\questiontitle{Transpose Convolution (stride 1)}
You are given an input $\mathbf{x}=\begin{pmatrix}0&1\\2&3\end{pmatrix}$ and a kernel
$\mathbf{w}=\begin{pmatrix}0&1\\2&3\end{pmatrix}$. What is the result of the transpose
convolution $\mathbf{w}\ast^{T}\mathbf{x}$ (stride 1)?
{\footnotesize
\begin{itemize}[leftmargin=1.8em]
  \choice $(42)$
  \choice $(14)$
  \choice $\begin{pmatrix}1&1&1&1\\2&2&2&2\\3&3&3&3\\4&4&4&4\end{pmatrix}$
  \choice $\begin{pmatrix}13&13&13\\13&13&13\\13&13&13\end{pmatrix}$
  \choice $\begin{pmatrix}0&2&1\\6&14&6\\2&3&0\end{pmatrix}$
  \choice $\begin{pmatrix}3&2\\1&0\end{pmatrix}$
  \choice $\begin{pmatrix}0&0&0&1\\0&0&2&3\\0&2&0&3\\4&6&6&9\end{pmatrix}$
  \choice $\begin{pmatrix}0&0&1\\0&4&6\\4&12&9\end{pmatrix}$
\end{itemize}}
\closequestion

\questiontitle{Transpose Convolution (stride 2)}
You are given an input $\mathbf{x}=\begin{pmatrix}1&-2\\-1&2\end{pmatrix}$ and a kernel
$\mathbf{w}=\begin{pmatrix}0&1\\2&3\end{pmatrix}$. What is the result of the transpose
convolution $\mathbf{w}\ast^{T}\mathbf{x}$ with \textbf{stride 2}?
{\footnotesize
\begin{itemize}[leftmargin=1.8em]
  \choice $\begin{pmatrix}0&1&-1\\2&-1&-1\\-4&-2&6\end{pmatrix}$
  \choice $\begin{pmatrix}3&2\\1&0\end{pmatrix}$
  \choice $\begin{pmatrix}0&1&0&-2\\2&3&-4&-6\\0&-1&0&2\\-2&-3&4&6\end{pmatrix}$
  \choice $\begin{pmatrix}0&2&1\\6&-14&6\\2&3&0\end{pmatrix}$
  \choice $(3)$
  \choice $(42)$
\end{itemize}}
\closequestion


% ==========================================================
\section{Autoencoders}
% ==========================================================

\questiontitle{Annotating the Generative-Model Taxonomy}
A figure shows three categories of generative models. Each row is a network with grey input/output
boxes and a latent variable $\mathbf{z}$ in the middle. Annotate the building blocks with
\emph{Generator}, \emph{Discriminator}, \emph{Encoder}, \emph{Decoder}, \emph{Flow}, \emph{Inverse Flow},
given that a rectangular block keeps the dimension constant while a trapezoidal block decreases or
increases it.
\examfigure[width=0.9\linewidth]{figures/exam-generative-model-taxonomy.png}
\begin{itemize}[leftmargin=1.8em]
  \item Top row: input $\to$ trapezoid $\to 0/1 \to \mathbf{z} \to$ trapezoid $\to$ output \\ \answerline{5cm}
  \item Middle row: input $\to$ trapezoid (down) $\to \mathbf{z} \to$ trapezoid (up) $\to$ output \\ \answerline{5cm}
  \item Bottom row: input $\to$ rectangle $\to \mathbf{z} \to$ rectangle $\to$ output \\ \answerline{5cm}
\end{itemize}
\closequestion

\questiontitle{Autoencoders (excluding VAEs)}
Which of the following statements about autoencoders (\textbf{not} including variational autoencoders)
are correct? (Select one or more.)
\begin{itemize}[leftmargin=1.8em]
  \choice The input--output Jacobian of an autoencoder is always a lower-triangular matrix.
  \choice Autoencoders do not map a single input to a single point in latent space, but to a distribution
        in latent space.
  \choice Autoencoders are always fully-connected networks.
  \choice All autoencoders use fewer neurons in the hidden layers than in the input layer.
  \choice Autoencoders consist of an encoder and a decoder, which are parametrized neural networks.
\end{itemize}
\closequestion

\questiontitle{Sparse and Denoising Autoencoders}
Which of the following statements about sparse / denoising autoencoders are correct?
(Select one or more.)
\begin{itemize}[leftmargin=1.8em]
  \choice Sparse AEs enforce sparsity on the representation via a regularization term.
  \choice Sparse AEs enforce hidden units to be inactive (i.e.\ activation $=0$) for most inputs.
  \choice A sparse AE can be trained on a copy task, where the input equals the target.
  \choice Sparse AEs force the \emph{input} units to be inactive (activation $=0$).
  \choice Sparse AEs enforce sparsity via early stopping.
  \choice A sparse AE cannot be trained on a denoising task (input noisy, target clean).
\end{itemize}
\closequestion

\questiontitle{Autoencoder Reconstruction Loss}
An autoencoder uses ReLU activations with encoder
$u=\mathrm{ReLU}\!\left(\begin{bmatrix}1&0&-1\\0&1&0\end{bmatrix}\mathbf{x}\right)$ and decoder
$\hat{\mathbf{x}}=\mathrm{ReLU}\!\left(\begin{bmatrix}1&0\\0&0\\1&0\end{bmatrix}u\right)$.
Compute the reconstruction loss $\loss(\mathbf{x},\hat{\mathbf{x}})=\sum_i (x_i-\hat x_i)^2$
for the input $\mathbf{x}=(1,1,2)\T$.

Answer: \answerline{3cm}
\closequestion


% ==========================================================
\section{Variational Autoencoders}
% ==========================================================

\questiontitle{Variational Autoencoders}
Which of the following statements about variational autoencoders (VAEs) are true? (Select one or more.)
\begin{itemize}[leftmargin=1.8em]
  \choice VAEs have an encoder--decoder structure, where encoder and decoder must be symmetric.
  \choice The prior distribution of the latent variables $\mathbf{z}$ must be a discrete distribution.
  \choice VAEs maximize a lower bound on the likelihood of the data.
  \choice VAEs automatically provide the marginal likelihood $p(\mathbf{x})$ at low computational cost.
  \choice VAEs cannot be used to generate images.
\end{itemize}
\closequestion

\questiontitle{VAE Objective Function}
Using the notation from the lecture, which of the following is the correct formulation of the
standard Variational Autoencoder (VAE) objective?
\begin{itemize}[leftmargin=1.8em]
  \choice $\loss(D,G)=\E_{\mathbf{x}\sim p_r}[\log D(\mathbf{x})]+\E_{\mathbf{z}\sim p_z}[\log(1-D(G(\mathbf{z})))]$
  \choice $\loss_{\theta,\phi}(\mathbf{x})=\E_{q_\phi(\mathbf{z}\given\mathbf{x})}\log p_\theta(\mathbf{x}\given\mathbf{z})-\KL{p_\theta(\mathbf{x})}{p_\theta(\mathbf{z})}$
  \choice $\loss(D,G)=\E_{\mathbf{x}\sim p_r}[\log D(\mathbf{x})]+\E_{\mathbf{z}\sim p_z}[D(G(\mathbf{z}))]$
  \choice $\loss_{\theta,\phi}(\mathbf{x})=\E_{q_\phi(\mathbf{z}\given\mathbf{x})}\log p_\theta(\mathbf{x}\given\mathbf{z})-\KL{q_\phi(\mathbf{z}\given\mathbf{x})}{p_\theta(\mathbf{z})}$
  \choice $\loss_\theta(\mathbf{x})=\KL{p_{\mathrm{data}}(\mathbf{x};\theta)}{p^{*}(\mathbf{x})}$
\end{itemize}
\closequestion

\questiontitle{Best Gaussian Approximation via KL}
\repeatnote\\
You have a fixed, multi-dimensional, non-Gaussian distribution $p(\mathbf{x})$ that you want to
approximate with a Gaussian $q(\mathbf{x})=\Normal(\boldsymbol\mu,\boldsymbol\Sigma)$, by minimizing
$\KL{p}{q}$ with respect to $\boldsymbol\mu$ and $\boldsymbol\Sigma$. Which expressions minimize the
KL-divergence? ($\E$ is expectation, $\Cov$ the covariance, $\mathbf{I}$ the identity matrix.)
\begin{itemize}[leftmargin=1.8em]
  \choice $\boldsymbol\mu=\E(\mathbf{x})$ and $\boldsymbol\Sigma=\Cov(\mathbf{x})$
  \choice $\boldsymbol\mu=\E(\mathbf{x})$ and $\boldsymbol\Sigma=\mathbf{I}$
  \choice $\boldsymbol\mu=\E(\mathbf{x}\T\mathbf{x})$ and $\boldsymbol\Sigma=\E(\mathbf{x}\T)$
  \choice $\boldsymbol\mu=\E(\mathbf{x}\T\mathbf{x})$ and $\boldsymbol\Sigma=\mathbf{I}$
  \choice $\boldsymbol\mu=\E(\mathbf{x})$ and $\boldsymbol\Sigma=\E(\mathbf{x}\mathbf{x}\T)$
\end{itemize}
\closequestion


% ==========================================================
\section{GANs and Wasserstein GANs}
% ==========================================================

\questiontitle{GAN Objective Function}
Using the lecture notation, which of the following is the correct formulation of the standard
objective of Generative Adversarial Networks (GANs)?
\begin{itemize}[leftmargin=1.8em]
  \choice $\min_G\max_D \loss(D,G)=\E_{\mathbf{x}\sim p_r}[\log D(\mathbf{x})]+\E_{\mathbf{z}\sim p_z}[\log(1-D(\mathbf{z}))]$
  \choice $\min_G\max_D \loss(D,G)=\E_{\mathbf{x}\sim p_r}[\log D(\mathbf{x})]+\E_{\mathbf{x}\sim p_r}[\log(1-D(G(\mathbf{x})))]$
  \choice $\min_G\max_D \loss(D,G)=\E_{\mathbf{x}\sim p_r}[\log G(\mathbf{x})]+\E_{\mathbf{z}\sim p_z}[\log(1-G(D(\mathbf{z})))]$
  \choice $\min_G\max_D \loss(D,G)=\E_{\mathbf{x}\sim p_r}[\log D(\mathbf{x})]+\E_{\mathbf{z}\sim p_z}[\log(1-D(G(\mathbf{z})))]$
  \choice $\min_G\max_D \loss(D,G)=\E_{\mathbf{x}\sim p_r}[\log D(\mathbf{x})]+\E_{\mathbf{z}\sim p_z}[D(G(\mathbf{z}))]$
  \choice $\min_G\max_D \loss(D,G)=\E_{\mathbf{x}\sim p_r}[\log D(\mathbf{x})]+\E_{\mathbf{z}\sim p_z}[\log(D(G(\mathbf{z})))]$
\end{itemize}
\closequestion

\questiontitle{Properties of a Trained GAN}
Consider a GAN that successfully produces images of apples. Which statements are correct?
(Select one or more.)
\begin{itemize}[leftmargin=1.8em]
  \choice The discriminator can be used to classify images as apple vs.\ non-apple.
  \choice The generator aims to learn the distribution of apple images.
  \choice The generator can produce unseen images of apples.
  \choice After training the GAN, the discriminator loss eventually reaches a constant value.
\end{itemize}
\closequestion

\questiontitle{Divergences and Distances}
Which of the following statements about divergences and distances are correct? (Select one or more.)
\begin{itemize}[leftmargin=1.8em]
  \choice JS and KL divergences show problems for distributions with disjoint supports --- the Wasserstein
        distance overcomes this.
  \choice A map $D:\mathcal{S}\times\mathcal{S}\to\R$ is a divergence iff it is non-negative and symmetric.
  \choice The Wasserstein-1 distance with Euclidean cost is a distance, not a divergence.
  \choice The Wasserstein distance is defined as a minimization over joint distributions (couplings),
        but this is infeasible in practice.
  \choice The Kantorovich--Rubinstein theorem lets us write the Wasserstein distance as a maximization
        over Lipschitz functions, which makes it tractable.
\end{itemize}
\closequestion

\questiontitle{KL, JSD and Wasserstein vs.\ Separation}
Let $p_1=\Normal(0,1)$ and $p_2=\Normal(x,1)$. As $x$ increases, you plot the
KL-divergence $\KL{p_1}{p_2}$, the Jensen--Shannon divergence $\mathrm{JSD}(p_1\|p_2)$, and the
Wasserstein-2 distance $\mathrm{WD}(p_1\|p_2)$ against $x$. One curve grows quadratically without
bound, one grows linearly, and one saturates at a finite value. Match each measure to its qualitative
behaviour.
\examfigure[width=0.86\linewidth]{figures/exam-kl-jsd-wasserstein-separation.png}
\begin{itemize}[leftmargin=1.8em]
  \item KL-divergence \answerline{5cm}
  \item Wasserstein-2 distance \answerline{5cm}
  \item Jensen--Shannon divergence \answerline{5cm}
\end{itemize}
\closequestion


% ==========================================================
\section{Metrics, Normalizing Flows \& Density Models}
% ==========================================================

\questiontitle{Inception Score and FID}
Tick the correct statements about the Inception Score (IS) and the Fréchet Inception Distance (FID).
(Select one or more.)
\begin{itemize}[leftmargin=1.8em]
  \choice FID can detect mode collapse, because mode collapse influences the covariance matrix of the
        generated data.
  \choice IS and FID are metrics that assess the quality of GAN-generated images.
  \choice Both IS and FID rely on features of the input images produced by AlexNet.
  \choice IS and FID can be used as GAN objectives that lead to faster training.
  \choice The computation of the FID is based on an explicit formula for the Wasserstein-2 distance
        between two Gaussian distributions.
\end{itemize}
\closequestion

\questiontitle{Normalizing Flows}
Which of the following statements about normalizing flows are correct? (Select one or more.)
\begin{itemize}[leftmargin=1.8em]
  \choice It is impossible to integrate a multi-layer perceptron as a transformation in a normalizing
        flow because MLPs are not invertible.
  \choice All multi-layer perceptrons with the same input and output dimension can be considered a
        normalizing flow.
  \choice Normalizing flows cannot contain functions with a lower-triangular Jacobian because they are
        not invertible.
  \choice Some normalizing flows use functions with a lower-triangular Jacobian to facilitate the
        computation of the determinant.
\end{itemize}
\closequestion

\questiontitle{Flow-Based Models}
Which of the following statements about flow-based models are correct? (Select one or more.)
\begin{itemize}[leftmargin=1.8em]
  \choice Theoretically, flow-based models can represent any distribution $p_x(\mathbf{x})$ if one starts
        from a sufficiently complex distribution; a simple base distribution $p_u(\mathbf{u})$ such as the
        uniform on the cube $(0,1)^D$ would not be sufficient.
  \choice Fitting flow-based models is often done with the KL-divergence.
  \choice Computing the Jacobi determinant of a neural network with $D$ inputs and $D$ outputs is in
        general of $\mathcal{O}(D^3)$.
  \choice For autoregressive flows the transformation $g_\phi$ is constructed so that its Jacobi
        determinant always has value 1.
\end{itemize}
\closequestion

\questiontitle{Density under a Linear Transform}
\repeatnote\\
Assume a two-dimensional vector $\mathbf{u}\sim p_u(\mathbf{u})$. Let
$T=\begin{pmatrix}2&0\\0&\tfrac12\end{pmatrix}$ and $\mathbf{x}=T\cdot\mathbf{u}$. Mark the correct
formula for the distribution $p_x(\mathbf{x})$.
{\footnotesize
\begin{itemize}[leftmargin=1.8em]
  \choice $p_u\!\left(\begin{pmatrix}\tfrac12&0\\0&2\end{pmatrix}\mathbf{x}\right)\cdot\tfrac12$
  \choice $p_u\!\left(\begin{pmatrix}\tfrac12&0\\0&2\end{pmatrix}\mathbf{x}\right)$
  \choice $p_u\!\left(\begin{pmatrix}\tfrac12&0\\0&1\end{pmatrix}\mathbf{x}\right)\cdot\tfrac12$
  \choice $p_u\!\left(\begin{pmatrix}\tfrac12&0\\0&1\end{pmatrix}\mathbf{x}\right)$
  \choice $p_u\!\left(\begin{pmatrix}2&0\\0&1\end{pmatrix}\mathbf{x}\right)\cdot 2$
  \choice $p_u\!\left(\begin{pmatrix}2&0\\0&1\end{pmatrix}\mathbf{x}\right)\cdot\tfrac12$
\end{itemize}}
\closequestion

\questiontitle{Density under a Linear Transform (Non-Diagonal)}
Let $\mathbf{u}\sim p_u(\mathbf{u})$ be a two-dimensional random vector and define
$\mathbf{x}=T\mathbf{u}$ with the non-diagonal matrix
$T=\begin{pmatrix}2&-2\\0&1\end{pmatrix}$. Derive the density $p_x(\mathbf{x})$ in terms of $p_u$,
stating $\det T$, $T^{-1}$, and the final formula.
\workingspace{4cm}
\closequestion

\questiontitle{Autoregressive Sequence Probability}
An autoregressive model $p_w(\mathbf{x})$ models sequences $(x_1,x_2,x_3)$ where each element is
$A$, $B$, or $C$:
\[
  p_w(x_1)=\tfrac13,\qquad
  p_w(x_2\given x_1)=\begin{cases}\tfrac12 & x_2=x_1\\[2pt]\tfrac14 & x_2\neq x_1\end{cases},\qquad
  p_w(x_3\given x_2,x_1)=\begin{cases}\tfrac12 & x_3=x_2\\[2pt]\tfrac14 & x_3\neq x_2\end{cases}.
\]
Compute $p_w(\mathbf{x})$ for $\mathbf{x}=(A,A,C)$, rounded to three decimals.

Answer: \answerline{3cm}
\closequestion


% ==========================================================
\section{Bayesian Deep Learning}
% ==========================================================

\questiontitle{Deep Ensembles and Marginalization}
The main difference between the frequentist and Bayesian approach is the marginalization over
parameters when defining the probability of a class label given input $\mathbf{x}$ and dataset
$\mathcal{D}$:
\[
  p(\mathbf{y}^{\mathrm{test}}\given\mathbf{x}^{\mathrm{test}},\mathcal{D})
  = \E_{\mathbf{w}\sim p(\mathbf{w}\given\mathcal{D})}\!\left[p(\mathbf{y}^{\mathrm{test}}\given\mathbf{w},\mathbf{x}^{\mathrm{test}})\right].
\]
Deep Ensembles approximate this full expectation with (weighted) averages. Which of the following
are valid approximations? (Select one or more.)
\begin{itemize}[leftmargin=1.8em]
  \choice $\tfrac1M\sum_{m=1}^{M} p(\mathbf{y}^{\mathrm{test}}\given\mathbf{w}_m,\mathbf{x}^{\mathrm{test}})$,
        where $\mathbf{w}_m$ are the weights of the $m$-th network of the ensemble.
  \choice $p(\mathbf{y}^{\mathrm{test}}\given\mathbf{w},\mathbf{x}^{\mathrm{test}})$ for a single trained network.
  \choice $\tfrac1M\sum_{m=1}^{M}\alpha_m\, p(\mathbf{y}^{\mathrm{test}}\given\mathbf{w}_m,\mathbf{x}^{\mathrm{test}})$
        with $\alpha_m=\softmax(\mathbf{z})_m$ for some random vector $\mathbf{z}$.
  \choice $\tfrac1M\sum_{m=1}^{M} p(\mathbf{y}^{\mathrm{test}}_m\given\mathbf{w},\mathbf{x}^{\mathrm{test}}_m)$,
        where $m$ denotes the $m$-th equally-sized fraction of the test data.
\end{itemize}
\closequestion


% ==========================================================
\section{Graph Neural Networks}
% ==========================================================

\questiontitle{Distinguishing Two Small Graphs}
Consider two graphs: $\mathcal{G}_1$ is a path of three nodes (a coloured centre node joined to two
leaves), and $\mathcal{G}_2$ is a star with a centre node joined to three leaves. Mark the correct
statements. (Select one or more.)
\examfigure[width=0.62\linewidth]{figures/exam-gnn-small-graphs.png}
\begin{itemize}[leftmargin=1.8em]
  \choice $\mathcal{G}_1$ and $\mathcal{G}_2$ can be distinguished by the Weisfeiler--Lehman
        (colour refinement) test.
  \choice All message-passing neural networks (MPNNs) can distinguish $\mathcal{G}_1$ and $\mathcal{G}_2$.
  \choice $\mathcal{G}_1$ and $\mathcal{G}_2$ are isomorphic graphs.
\end{itemize}
\closequestion

\questiontitle{Weisfeiler--Lehman Necessary Condition}
Two graphs $A$ and $B$ are given. Do they fulfil the necessary condition for graph isomorphism
given by the Weisfeiler--Lehman test (i.e.\ do they produce the same stable colour histogram)?
\begin{itemize}[leftmargin=1.8em]
  \choice True
  \choice False
\end{itemize}
\closequestion

\questiontitle{Matching Graphs to Adjacency Matrices}
Four graphs on six nodes are shown: $A$ is a plain $6$-cycle (hexagon), $B$ is two triangles,
$C$ is another $6$-cycle drawn differently, and $D$ is a $6$-cycle with one extra chord.
For each adjacency matrix, name the matching graph ($A$, $B$, $C$, $D$, or none).
\examfigure[width=0.94\linewidth]{figures/exam-gnn-candidate-graphs.png}
\begin{itemize}[leftmargin=1.8em]
  \item All six row-sums equal $2$ and the edges form a single closed cycle. \hfill\answerline{3cm}
  \item Two row-sums equal $1$ and the rest equal $2$ (an open path, not a cycle). \hfill\answerline{3cm}
  \item All six row-sums equal $2$ but the edges split into two separate $3$-cycles. \hfill\answerline{3cm}
  \item Four row-sums equal $2$ and two equal $3$ (a cycle plus one chord). \hfill\answerline{3cm}
\end{itemize}
\closequestion

\questiontitle{Group Action in Graph NNs}
Which of the following is a common group action used in graph neural networks?
\begin{itemize}[leftmargin=1.8em]
  \choice Rotation
  \choice Scaling
  \choice Translation (permutation / relabelling of nodes)
  \choice Cropping
  \choice Clustering
\end{itemize}
\closequestion

\questiontitle{Equivariance of an Attention-Pooling Layer}
A layer implements $\mathbf{h}^{(l+1)}=M\,\underbrace{\softmax(\beta\,M\T\mathbf{h}^{(l)})}_{=:\,\mathbf{p}}$,
where $\mathbf{h}^{(l)}\in\R^{d}$, $\beta>0$ is a hyperparameter, and $M\in\R^{d\times J}$ holds $J$
pattern columns. Which statements about this layer are true? (Select one or more.)
\begin{itemize}[leftmargin=1.8em]
  \choice $\mathbf{h}^{(l+1)}$ is equivariant to permutations of the columns of $M$.
  \choice $\mathbf{p}$ is invariant to permutations of the columns of $M$.
  \choice $\mathbf{p}$ is equivariant to permutations of the columns of $M$.
  \choice $\mathbf{h}^{(l+1)}$ is invariant to permutations of the columns of $M$.
  \choice $\mathbf{p}$ is equivariant to permutations of the rows of $M$.
  \choice This layer cannot be used in deep learning because the gradient
        $\partial\mathbf{h}^{(l+1)}/\partial\mathbf{h}^{(l)}$ cannot be backpropagated.
\end{itemize}
\closequestion


% ==========================================================
\section{Theory \& Foundations}
% ==========================================================

\questiontitle{Universal Approximator Theorem}
Mark the correct statements about the Universal Approximator Theorem (UAT) for neural networks.
(Select one or more.)
\begin{itemize}[leftmargin=1.8em]
  \choice For a real-valued, continuous function $h$ on $[-\pi,\pi]^2$, the existence of an approximating
        network can be derived with the help of Fourier series.
  \choice The UAT is a theoretical statement about the expressive power of neural networks.
  \choice The UAT gives a useful algorithm to compute the optimal weights of a network that should
        approximate a given continuous function.
  \choice The algorithm related to the UAT is easy to understand and implement, but very costly and
        therefore hardly used in practice.
  \choice The UAT shows that any function can be approximated by any given network with arbitrary
        precision; there are no restrictions on the function or on the network structure.
\end{itemize}
\closequestion

\questiontitle{Newton's Method vs.\ Gradient Direction}
On a quadratic error surface (elliptical level curves), two arrows are drawn from a point: one points
straight at the minimum at the centre of the ellipses, the other points perpendicular to the local
level curve. Label each arrow as \emph{Gradient} or \emph{Newton's Method}.
\examfigure[width=0.9\linewidth]{figures/exam-newton-gradient-directions.png}
\begin{itemize}[leftmargin=1.8em]
  \item Arrow pointing directly to the centre (minimum) \\ \answerline{4cm}
  \item Arrow perpendicular to the level curve (steepest descent) \\ \answerline{4cm}
\end{itemize}
\closequestion


% ==========================================================
%  ANSWER KEY
% ==========================================================
\newpage
\section*{Answer Key}
\addcontentsline{toc}{section}{Answer Key}

\begin{enumerate}[label=\textbf{Q\arabic*.}, leftmargin=3em, itemsep=4pt]
  \item \textbf{c, e.} AlexNet \emph{does} have fully-connected layers (a wrong); it uses both
        strided convolutions \emph{and} max-pooling (b wrong); it uses ReLU, not BatchNorm, which came
        later (d wrong).
  \item \textbf{a, d.} VGG uses stacked $3\times3$ convolutions (two $3\times3$ replace one $5\times5$
        at the same receptive field, so b is wrong); it is a plain classifier, not an encoder--decoder
        (c wrong); most parameters sit in the large FC layers.
  \item \textbf{c, d.} AlexNet used Local Response Normalization, not BatchNorm (a wrong); LayerNorm is
        mostly used in transformers / sequence models, while CNN classifiers still use BatchNorm (b wrong).
  \item \textbf{a.} Normalization is central to training deep ResNets (b wrong); pre-trained ResNets are
        widely used for transfer learning (c wrong); residual blocks also work with FC layers (d wrong);
        leaky ReLUs can be used (e wrong).
  \item \textbf{b.} A $1\times1$ conv with 16 kernels keeps the spatial size and sets the channel depth to
        16: output $[128,128,16]$. Stride 2 or $2\times2$ pooling would shrink the spatial resolution.
  \item Increase stride: \textbf{no increase} (only the last layer changes, so its stride does not enlarge
        each neuron's RF). Dilated convolution: \textbf{strong increase}. Number of kernels:
        \textbf{no increase}. Kernel size: \textbf{small increase}.
  \item \textbf{c} ($0.5$). With all weights/biases zero, every pre-activation is $0$ and
        $\sigma(0)=0.5$ at each layer.
  \item \textbf{b} (may be positive or negative). The hidden activations are all the identical positive
        value $0.5$, so the rows of $\mathbf{W}^{(2)}$ receive equal updates, but their sign is set by the
        per-output error $(\hat y-y)$, which can be positive or negative. (They are not all zero, so d is wrong.)
  \item \textbf{a.} Pixel-wise labels are needed. This is semantic segmentation, not object detection
        (b wrong); CNNs with a (cross-entropy) loss are appropriate, not GANs (c, d wrong); IoU/mIoU is a
        standard segmentation metric (e wrong).
  \item $\mathrm{IoU}=\dfrac{|\text{intersection}|}{|\text{union}|}=\dfrac{5}{10+8-5}=\dfrac{5}{13}\approx\mathbf{0.385}$.
  \item \textbf{a, b, d, e.} Faster R-CNN replaces selective search with an RPN; RoI-pooling/align act on
        feature maps; detection is a multi-task regression + classification problem; Fast R-CNN has a
        regression head and a classification head. (c is debatable/considered false: R-CNN runs a CNN on
        each warped proposal but is not characterised by ``AlexNet'' in the lecture.)
  \item Faster R-CNN \textbf{two}-stage; YOLO \textbf{one}-stage; Mask R-CNN \textbf{two}-stage;
        Fast R-CNN \textbf{two}-stage; R-CNN \textbf{two}-stage; SSD \textbf{one}-stage.
  \item \textbf{h} $=\begin{pmatrix}0&0&1\\0&4&6\\4&12&9\end{pmatrix}$ (full stride-1 transpose convolution
        of two $2\times2$ inputs gives a $3\times3$ output).
  \item \textbf{c} $=\begin{pmatrix}0&1&0&-2\\2&3&-4&-6\\0&-1&0&2\\-2&-3&4&6\end{pmatrix}$ (stride-2
        transpose convolution spaces the kernel copies two apart, giving a $4\times4$ output).
  \item Top row = GAN: trapezoid in $=$ \textbf{Discriminator}, trapezoid out $=$ \textbf{Generator}.
        Middle row = (V)AE: down $=$ \textbf{Encoder} $q(\mathbf{z}\given\mathbf{x})$, up $=$
        \textbf{Decoder} $p(\mathbf{x}\given\mathbf{z})$. Bottom row = flow: first rectangle $=$
        \textbf{Flow}, second $=$ \textbf{Inverse Flow}.
  \item \textbf{e.} Only ``encoder + decoder are parametrized networks'' always holds. The lower-triangular
        Jacobian (a) is not generally true; mapping to a \emph{distribution} (b) is the VAE; AEs need not be
        fully-connected (c) nor undercomplete (d) --- sparse / overcomplete AEs exist.
  \item \textbf{a, b, c.} Sparsity is enforced by a regularization term that keeps most \emph{hidden} units
        inactive, and such an AE can still be trained on a copy task. It does not silence \emph{input} units
        (d wrong), does not rely on early stopping (e wrong), and sparsity and denoising can be combined (f wrong).
  \item \textbf{$6$.} $u=\mathrm{ReLU}([\,1-2,\,1\,])=\mathrm{ReLU}([-1,1])=[0,1]$;
        $\hat{\mathbf{x}}=\mathrm{ReLU}([0,0,0])=[0,0,0]$; loss $=(1{-}0)^2+(1{-}0)^2+(2{-}0)^2=1+1+4=6$.
        (If a $\tfrac1n$ MSE is used instead of the sum, the value is $2.0$.)
  \item \textbf{c.} VAEs maximize the ELBO, a lower bound on the data log-likelihood. Encoder/decoder need
        not be symmetric (a); the prior is typically a continuous (multivariate Normal) distribution, not
        discrete (b); they do not give $p(\mathbf{x})$ cheaply (d); they can generate images (e).
  \item \textbf{d.} The ELBO is
        $\E_{q_\phi(\mathbf{z}\given\mathbf{x})}\log p_\theta(\mathbf{x}\given\mathbf{z})
        -\KL{q_\phi(\mathbf{z}\given\mathbf{x})}{p_\theta(\mathbf{z})}$. Option b has the wrong KL arguments.
  \item \textbf{a.} The KL-optimal Gaussian matches the first two moments: $\boldsymbol\mu=\E(\mathbf{x})$
        and $\boldsymbol\Sigma=\Cov(\mathbf{x})$. Option e uses the second moment $\E(\mathbf{x}\mathbf{x}\T)$,
        which overestimates the covariance (it omits $-\E(\mathbf{x})\E(\mathbf{x})\T$).
  \item \textbf{d.} The standard GAN value function is
        $\E_{\mathbf{x}\sim p_r}[\log D(\mathbf{x})]+\E_{\mathbf{z}\sim p_z}[\log(1-D(G(\mathbf{z})))]$.
  \item \textbf{b, c.} The generator learns the data distribution and can produce unseen samples. The
        discriminator only learns real-vs-fake, not apple-vs-non-apple (a wrong); the discriminator loss
        reaching a constant is not guaranteed in general (d wrong).
  \item \textbf{a, c, d, e.} A divergence must be non-negative and zero iff the distributions are equal, but
        \emph{need not be symmetric} --- so b (non-negative + symmetric) is the definition of a distance, not
        a divergence.
  \item KL-divergence: grows quadratically without bound (\textbf{purple, solid}). Wasserstein-2: grows
        \emph{linearly} (\textbf{green, long-dashed}). Jensen--Shannon: \emph{saturates} at a finite value
        (\textbf{orange, short-dashed}).
  \item \textbf{a, b, e.} FID compares the mean and covariance of real vs.\ generated features (so it detects
        mode collapse via the covariance) and is in fact the closed-form Wasserstein-2 distance between two
        Gaussians; IS and FID assess GAN image quality. They use Inception features, not AlexNet (c wrong),
        and are not training objectives (d wrong).
  \item \textbf{d.} A lower-triangular Jacobian makes the determinant the product of the diagonal, which is
        cheap. Standard MLPs are generally not invertible, but it is not ``impossible'' to design invertible
        ones (a wrong); not all equal-dimension MLPs are flows (b); lower-triangular maps can be invertible (c).
  \item \textbf{b, c.} Flows are fit by maximum likelihood, i.e.\ minimizing a KL-divergence; the general
        Jacobi determinant costs $\mathcal{O}(D^3)$. A uniform base \emph{is} sufficient, so a is wrong;
        an autoregressive flow's Jacobian is triangular but its determinant is \emph{not} always $1$ (d wrong).
  \item \textbf{b.} $p_x(\mathbf{x})=\dfrac{1}{|\det T|}\,p_u(T^{-1}\mathbf{x})$ with
        $T^{-1}=\begin{pmatrix}\tfrac12&0\\0&2\end{pmatrix}$ and $\det T=2\cdot\tfrac12=1$, so the prefactor is $1$.
  \item $\det T=2\cdot1-(-2)\cdot0=2$, so $T^{-1}=\tfrac12\begin{pmatrix}1&2\\0&2\end{pmatrix}=\begin{pmatrix}\tfrac12&1\\0&1\end{pmatrix}$, giving
        $p_x(\mathbf{x})=\dfrac{1}{|\det T|}\,p_u(T^{-1}\mathbf{x})=\tfrac12\,p_u\!\left(\begin{pmatrix}\tfrac12&1\\0&1\end{pmatrix}\mathbf{x}\right)$.
  \item $p_w(A,A,C)=\tfrac13\cdot\tfrac12\cdot\tfrac14=\tfrac{1}{24}\approx\mathbf{0.042}$
        ($x_2=A=x_1\Rightarrow\tfrac12$; $x_3=C\neq A=x_2\Rightarrow\tfrac14$).
  \item \textbf{a, c.} A uniform average over the ensemble members (a) and a weighted average (c) both
        approximate the posterior expectation. A single network (b) is not a Bayesian average; splitting the
        \emph{test data} into fractions (d) is not ensembling.
  \item \textbf{a.} They have different structure (in fact different node counts / degree sequences), so the
        WL test distinguishes them. Not all MPNNs can (expressivity depends on the aggregation), so b is wrong;
        they are not isomorphic (c wrong).
  \item \textbf{True.} The two graphs produce the same stable WL colour histogram, so the necessary
        condition for isomorphism is fulfilled.
  \item Single $6$-cycle (all degrees $2$, one cycle) $\to$ $A$ (and the differently-drawn cycle $\to C$);
        open path (two degree-$1$ nodes) $\to$ \emph{none of the above}; two separate $3$-cycles $\to$ $B$;
        cycle with one chord (two degree-$3$ nodes) $\to$ $D$. Match each matrix by its row-sums (node degrees)
        and whether the edges form one cycle, two triangles, or a cycle-plus-chord.
  \item \textbf{c} (Translation / permutation). GNNs are permutation invariant/equivariant; rotation, scaling,
        cropping, and clustering are not the defining group action here.
  \item \textbf{c, d.} Permuting the columns of $M$ permutes the rows of $M\T$ and hence the entries of
        $\mathbf{p}$ (so $\mathbf{p}$ is \emph{equivariant} to column permutations, c). Then
        $\mathbf{h}^{(l+1)}=M\mathbf{p}$ is unchanged because the column reorder of $M$ cancels the row reorder
        of $\mathbf{p}$ (so $\mathbf{h}^{(l+1)}$ is \emph{invariant}, d). The layer is differentiable (f wrong).
  \item \textbf{a, b.} The UAT can be argued via Fourier series and is a statement about expressive power.
        It does \emph{not} give an algorithm for the weights (c, d wrong), and it does have conditions
        (the function is continuous on a compact set and the network needs at least one hidden layer), so
        e is wrong.
  \item Arrow to the centre $=$ \textbf{Newton's Method} (uses curvature to point straight at the minimum);
        arrow perpendicular to the level curve $=$ \textbf{Gradient} (steepest-descent direction).
\end{enumerate}


\end{document}
